% VI.04.P8
\subsection*{Problem VI.04.P8 --- Equalizers and fixed-point loci}
\label{prob:vi04-p8}
\paragraph{Problem.}
Let \(f,g:X\to Y\subseteq\mathbb A^m\) be morphisms.  Prove that
\[
E=\set{x\in X:f(x)=g(x)}
\]
is Zariski closed in \(X\).  If \(f_i,g_i\in k[X]\) are their coordinate functions, show
\[
E=V_X(f_1-g_1,\dots,f_m-g_m).
\]
Deduce that the fixed-point set of every endomorphism \(X\to X\) is closed.
\ifdefined\FullProblemDossiers
\paragraph{Why this problem matters.}  Equality of morphisms is algebraic because it is
equality of coordinate functions.  This simple fact underlies diagonals, graphs, and many
later incidence constructions.
\paragraph{Guided hints.}  Two points of \(Y\subseteq\mathbb A^m\) are equal iff all their
ambient coordinates agree.
\paragraph{Solution.}
For \(x\in X\), \(f(x)=g(x)\) iff \(f_i(x)=g_i(x)\) for every \(i\).  Thus
\[
E=V_X(f_1-g_1,\dots,f_m-g_m),
\]
which is closed.  Taking \(g=\operatorname{id}_X\) gives the fixed-point locus of an
endomorphism.
\paragraph{What happened mathematically.}  The diagonal relation in affine space is cut out
by coordinate differences, so equalizers of affine morphisms remain algebraic.
\paragraph{Extensions.}  Reinterpret \(E\) as the inverse image of the diagonal in
\(Y\times Y\) under \((f,g)\).
\paragraph{Provenance.}  New synthesis problem preparing the later scheme-theoretic diagonal.
\fi
